\documentclass[10pt,a4paper]{amsart}
\usepackage[margin=19mm]{geometry}
\usepackage{amsmath,amssymb,mathtools}
\usepackage[scr=rsfs,cal=euler]{mathalfa}
\usepackage{xfrac}
\usepackage[unicode,hidelinks,bookmarks=false]{hyperref}
\newcommand{\ie}{i.e.\ }
\newcommand{\cf}{cf.\ }
\newcommand{\resp}{respectively\ }
\newcommand{\Subsection}{\subsection}
\providecommand{\nobreakdash}{\nobreak\hbox{-}}
\setlength{\emergencystretch}{2em}
\multlinegap0pt
\usepackage{tikz}
\usetikzlibrary{cd}
\usepackage{mathtools}
\usepackage{amssymb}
\usepackage{xcolor}
\newtheorem{definition}{Definition}
\newtheorem{corollary}{Corollary}
\newtheorem{theorem}{Theorem}
\newcommand{\Ab}{{\mathsf{Ab}}}
\newcommand{\N}{{\mathbb N}}
\newcommand{\Rings}{{\mathsf{Rings}}}
\newcommand{\Union}{\bigcup}
\newcommand{\Z}{{\mathbb Z}}
\newcommand{\ab}[1]{\langle {#1} \rangle}
\newcommand{\sat}{{\rm sat}}
\renewcommand{\tilde}{\widetilde}
\title{A universal construction of $p$-typical Witt vectors of associative rings}
\author{Supriya Pisolkar, Biswanath Samanta}
\date{Source: arXiv:2601.20536v1 (CC BY 4.0)}
\begin{document}
\maketitle

\noindent\emph{Source and licence.} A universal construction of $p$-typical Witt vectors of associative rings. Version arXiv:2601.20536v1 (CC BY 4.0), distributed by arXiv under the Creative Commons Attribution 4.0 International licence, http://creativecommons.org/licenses/by/4.0/, as recorded in the arXiv licence metadata for this exact version and shown on the abstract page https://arxiv.org/abs/2601.20536v1. Full licence notice: \texttt{licenses/arxiv-2601.20536-CC-BY-4.0.txt}.\par\smallskip

\noindent\textit{Editorial scope.} Counted unit: the article's theorem universalweakprewitt (source0.tex lines 557--558). The article's complete original proof of it is delivered with it (source0.tex lines 560--636, one \texttt{proof} environment of 77 lines). No mathematics is written, repaired or re-derived: the statement and the proof are the article's own lines, in the article's own notation, and no retained line is substituted or re-flowed.  Retained uncounted as the source-local context the proof needs: lines 246--259; lines 338--341; lines 349--381; lines 512--527; lines 539--541; lines 549--552.  Nothing was omitted from any retained span, so the package carries no omission record.  No retained line contains a citation, so the wrapper carries no bibliography and no bibliography entry.  The retained lines contain the article's own diagram code, which the wrapper typesets with the standard tikz packages; no image file is delivered and the package declares no asset dependency.  Editorial formatting only, in addition: an AMS \texttt{amsart} preamble that restates the article's own declarations and macros used by the retained lines, namely the environments \texttt{definition}, \texttt{corollary}, \texttt{theorem} on separate counters, together with the standard packages those lines need.  The retained environments print in this document as Definition 1, Corollary 1, Theorem 1 in order of appearance, whereas the article prints the same items with the numbering of its own full document.  The complete native source of the article as distributed in the arXiv e-print for this exact version is bundled unchanged as \texttt{sources/193528/source0.tex}.
\begin{definition}[pre-Witt functor]\label{def:prewitt}
    A pre-Witt functor is a functor $F:\Rings \to \Ab$ which has functorial maps 
\begin{enumerate}
    \item Verschiebung operator: $V: F(R)\to F(R)$ (group homomorphism)
    \item Teichm\"{u}ller map : $\ab{\ }: R\to F(R)$ (a set map)
\end{enumerate}
which satisfy the following properties 
\begin{enumerate}\label{properties}
    \item $\ab{0}=0$. If $p\neq 2$ then $\ab{-x} = -\ab{x}$ for all $x \in R$.
    \item $x \mapsto V\ab{x^p} -p\ab{x}$ is an additive map from $R \to F(R)$. 
    \item $F(R)$ is complete w.r.t to the filtration $\{V^nF(R)\ | \ n \in \N_0 \}$.
    \item $R$ and $\bar{R}$ are $p$-torsion free $\implies F(R)$ is $p$-torsion free, where $\bar{R} := \frac{R}{\ab{[R,R]}}$ and  $\ab{[R,R]}$ is the ideal generated by the commutator subgroup $[R,R]$. 
\end{enumerate}
\end{definition}

\begin{definition}[$X_I$ and $X_I^{\sat}$]\label{def:XI}
For $R \in \Rings$, consider the free presentation $$0 \to I \to A \to R \to 0 $$ where $A:=\Z \{ R\}$ be the non-commutative polynomial ring on the set $R$.  Let $X_I(A) \subset X(A)$ be the closed subgroup generated by $$\{V^n\ab{a}-V^n\ab{b} \ | \ n\geq 0,  \ a,b\in A \ \text{ such that } a-b\in I\}.$$ We will denote $X_I$ as a short for $X_I(A)$. We denote the $p$-saturation of $X_I$ by 
$$X_I^{\sat} := \{ \alpha \in X(A) \mid \ p^{\ell}\alpha \in X_I  \text{ for some } \ell \in \N_0 \} $$  
\end{definition}

\begin{definition}[$E(R)$]\label{def:e}
For $R \in \Rings$ we denote by $\bar{R}$ the quotient ring $\frac{R}{\ab{[R,R]}}$ where $\ab{[R,R]}$ is the ideal generated by the commutator subgroup $[R,R]$. Suppose there exists a morphism $\phi: S \to R$ in $\Rings$ where $S$, $\bar{S}$ are $p$-torsion free. We consider the free presentations of $R$ and $S$ as follows. 

\[
\begin{tikzcd}[column sep=small, row sep=1.5em]
0 \arrow[r] & J \arrow[d, dashed] \arrow[r] & \Z\{S\} \arrow[d, dashed,"\tilde{\phi}"] \arrow[r, "\pi_2"] & S \arrow[d, "\phi"] \arrow[r] & 0 \\
0 \arrow[r] & I \arrow[r] & \Z\{R\} \arrow[r, "\pi_1"] & R \arrow[r] & 0
\end{tikzcd}
\tag{1}
\] 

Here, $\tilde{\phi}$ is a lift of $\phi$. Applying the functor $X$ we get the following diagram.
\[
\begin{tikzcd}[column sep=small, row sep=1.5em]
 & X_J \arrow[d, dashed] \arrow[r] & X(\Z\{S\}) \arrow[d, dashed,"\tilde{\phi}"] \arrow[r] & X(S) \arrow[d, "\phi"] \arrow[r] & 0 \\
 & X_I \arrow[r] & X(\Z\{R\}) \arrow[r] & X(R) \arrow[r] & 0
\end{tikzcd}
\tag{2}
\] 

Following the notation as in the Definition \ref{def:XI}, we define 
$$\tilde{X_I} : = \text{ the closed subgroup of } X(\Z\{R\}) \text{ generated by  } \{X_I, \Sigma_I\}$$
where for a lift $\tilde{\phi}$ as above
\[
\begin{aligned}
\Sigma_I \;:=\; \bigl\{\, \alpha \in X(\Z\{R\}) \ \bigm|\;
&  \alpha = \tilde{\phi}(\beta) \text{ for some } \beta \in X_J^{\mathrm{sat}}
\bigr\}.
\end{aligned}
\]
We leave it to the reader to check that the set $\Sigma_I$ is independent of the lift $\tilde{\phi}$.
Finally, we define $$E(R):= \frac{X(\Z\{R\})}{\tilde{X_I}}$$ 
\end{definition}

\begin{definition}[The functor $C$]\label{def:C}
For any associative ring $R$ and a prime $p\neq 2$, we first define a group $G(R)$ as follows. Let 
\begin{enumerate}
\item $\tilde{G}(R)$ be the free abelian group generated by symbols $\{ V^n_r \ | \ n\in \N_0 , \ 0\neq r\in R\}$. Define $V^n_0:=0$ for all $n\geq0$. 
\item A set map $R \xrightarrow{\ab{\ } } \tilde{G}(R)$ given by $\ab{r} := V^0_r$.
\item A homomorphism $\tilde{G}(R) \xrightarrow{V} \tilde{G}(R)$ defined by $V(V^n_r):= V^{n+1}_r$.
\item $H(R)\subset \tilde{G}(R)$ be the subgroup generated by the set: 
$$  \Union_n\left(\{(V^n_{(x+y)^p} - p V^{n-1}_{x+y}) - (V^n_{x^p} - p V^{n-1}_{x}) - (V^n_{y^p} - p V^{n-1}_{y})\ | \ x, y \in R \} \Union \{V^n_r+V^n_{-r} \ | \ r \in R\}\right)$$ 
\item $\tilde{H}(R)$ be the $p$-saturation of $H(R)$, i.e.
$ \tilde{H}(R):= \{ \alpha \in \tilde{G}(R) \ | \ p^\ell \alpha \in H(R) \ \text{for some } \ell >0\}$.
\item Let $G^0(R)$ denote the completion of $\tilde{G}(R)/\tilde{H}(R)$ by the $V$-filtration. Let $G(R)$ be the quotient of $G^0(R)$ modulo the closed subgroup generated by $p$-power torsion elements. Note that the construction of $G(R)$ is functorial in $R$.\\
\end{enumerate}
Consider the free presentation $0 \to I \to A \to R \to 0 $ where $A:=\Z\{R\}$. For $x \in R$ we denote the corresponding variable in $A$ by $[x]$. Let $G_I(A) \subset G(A)$ is the closed subgroup generated by $$\{V^n_a-V^n_b \ | \ n\geq 0,  \ a,b\in A \ \text{ such that } a-b\in I\}.$$

\noindent Now define $$C(R):= \frac{G(A)}{\tilde{G_I}}$$ where $\tilde{G_I}$ is subgroup generated by $G_I$ and $\Sigma_I$ where $\Sigma_I$ and $G_I$ are defined similar to the definition of $\Sigma_I$ and $X_I$ as in Definition \ref{def:e}. 
\end{definition} 

\begin{corollary}\label{ceqg}
If $A=\Z\{S\}$ is a free algebra, then $C(A) \cong G(A)$.
\end{corollary}

\begin{definition}[weak pre-Witt functor]\label{def:weakpre} A functor $F: \Rings \to  \mathsf{Ab}$ is said to be a weak pre-Witt functor if it satisfies the properties (1), (2) and (3) of \ref{properties} and the following property, which is a weaker version of  \ref{properties}(4): \\

$(4')$ $A$ is a free polynomial ring in $\Rings$ then $F(A)$ is $p$-torsion free
\end{definition} 

\begin{theorem}\label{universalweakprewitt} Let $p\neq 2$.  The functor $C: \Rings \to \Ab$ as defined in \ref{def:C}, is a weak pre-Witt functor.  Moreover, given any pre-Witt functor $F$, there exists a natural transformation $C \to F$. In particular, there is a natural transformation $C\xrightarrow{\eta} E$ which is compatible with $V, \ab{ \ }$.
\end{theorem}

\begin{proof} It follows from the construction of $C$ that, $C$ satisfies properties (1), (2) and (3) of Definition \ref{properties}. The Corollary \ref{ceqg} implies that $C$ also satisfies the property $(4')$ of the Definition \ref{def:weakpre}. Hence $C$ is a weak pre-Witt functor. We will now prove that given any pre-Witt functor $F$ there exists a natural transformation $\eta: C \to F$. The idea is to first establish a group homomorphism in the case of a free algebra and then extend the result to any $R \in \Rings$. We will prove this in the following steps. \\

\noindent Step 1: We will show that, for a free non-commutative polynomial ring $A$ there exists a group homomorphism $C(A) \to F(A)$. Consider the presentation $$0 \to 0 \to A \xrightarrow{id} A \to 0$$ We define the homomorphism, 

$$\eta : \tilde{G}(A) \to F(A) \ \text{given by} \  V^n_{a}  \mapsto V^n\ab{a} \ \forall \  a \in A, \ n \geq 0.$$ 

\vspace{1mm}
\noindent Since $F$ satisfies the properties (1)-(2) of Definition \ref{def:weakpre}, $\eta(\tilde{H}(A))=0$ in $F(A)$. So, we get a homomorphism,

 $$\tilde{\eta}:\tilde{G}(A)/\tilde{H}(A) \to F(A).$$ 
 
\vspace{1mm}
\noindent As $\eta$ is compatible with $V$, taking $V$-completion, we get the induced group homomorphism, denoted as $\eta_A: G(A) \to F(A)$. \\

\noindent{Step 2:} To show that for $R \in \Rings$, there exists a group homomorphism $C(R) \to F(R)$.\\

\noindent Consider the presentation $0 \to I \to \Z\{R\} \stackrel{\pi_1}{\longrightarrow} R \to 0$ of $R$. By Step (1), we have a group homomorphism $F(\pi_1) \circ \eta_{\Z\{R\}} : G(\Z\{R\}) \to F(R)$. \\
\vspace{2mm}

\noindent It is easy to see that for $a, b \in \Z\{R\}$ such that $a-b \in I$, $G(\pi_1)(V^n_a-V^n_b)= V^n\ab{\pi_1(a)}-V^n\ab{\pi_1(b)} = 0$, thus $F(\pi_1)\circ\eta_{\Z\{R\}}(G_I) = 0$. Hence there exists a group homomorphism $\eta: G(\Z\{R\})/G_I \to F(R)$. It remains to show that $F(\pi_1) \circ \eta_{\Z\{R\}}(\tilde{G_I}) = 0$ \ \\

\noindent Let $S \in \Rings$ such that $S, \bar{S}$ are $p$-torsion free and $\phi: S \to R$ be a ring homomorphism. We now recall here the definition of $\Sigma$ and the diagram (1) as in the given in Definition \ref{def:e}. 
\[
\begin{tikzcd}[column sep=1.5em, row sep=1.5em]
0 \arrow[r] & J \arrow[d, dashed] \arrow[r] & \Z\{S\} \arrow[d, dashed,"\tilde{\phi}"] \arrow[r, "\pi_2"] & S \arrow[d, "\phi"] \arrow[r] & 0 \\
0 \arrow[r] & I \arrow[r] & \Z\{R\} \arrow[r, "\pi_1"] & R \arrow[r] & 0
\end{tikzcd}
\tag{1}
\] 
\noindent Consider the map $f: G(\Z\{S\}) \to G(\Z\{S\})/G_J$. Then 
$$\Sigma := \{ \alpha \in G(\Z\{R\}) \mid \alpha = G(\tilde{\phi})(\beta), \ \beta \ \text{is a p-torsion in \ }  G(\Z\{S\})/G_J \}$$

\noindent We have the following diagram 
\begin{center}
\begin{tikzcd}
&
G(\Z\{S \})
\ar{dl}[swap, sloped, near start]{\eta_{\Z\{S\}}}
\ar{rr}{f}
\ar[]{dd}[near start]{\tiny{G(\tilde{\phi})}}
& & G(\Z\{S\})/G_J 
\ar{dd}{}
\ar{dl}[swap, sloped, near start]{g}
\\
F(\Z\{S\}) 
\ar[crossing over]{rr}[swap, near start]{F(\pi_2)}
\ar{dd}[swap]{F(\tilde{\phi})}
& & F(S)
\ar[]{dd}[near start]{F(\phi)}
\\
&
G(\Z\{R\})
\ar[near start]{rr}{}
\ar[sloped]{dl}{\eta_{\Z\{R\}}}
& &  G(\Z\{R\})/G_I
\arrow[dl, densely dotted]
\\
F(\Z\{R\})
\ar[rr, "F(\pi_1)"]
& & F(R)
\ar[crossing over, leftarrow, near start]{uu}{}
\end{tikzcd}
\end{center}

\noindent Let $\alpha \in \Sigma$. Then $\alpha = G(\tilde{\phi})(\beta)$ for some $\beta \in G_J^{\sat}$. Thus $p^{\ell}\beta \in G_J$ for some $\ell \geq 0$. In other words $f(\beta)$ is a $p$-torsion element in $G(\Z\{S\})/G_J$. Since $\eta_{\Z\{S\}}(G_J) \subseteq F_J$, we have the induced group homomorphism denoted by $g: G(\Z\{S\})/G_J \to F(S)$, which takes $f(\beta)$ to a $p$-torsion element in $F(S)$. Since $F$ is a pre-Witt functor and $S, \bar{S}$ are $p$-torsion free,  $F(S)$ is $p$-torsion free. Thus $g \circ f(\beta) = 0$. Thus,
$$(F(\phi) \circ g \circ f)(\beta) = (F(\phi)\circ F(\pi_2)\circ \eta_{\Z\{S\}})(\beta) = 0$$
Since $F$ is functor, we have the commutativity of the front square and thus,
$$(F(\phi) \circ F(\pi_2) \circ \eta_{\Z\{S\}})(\beta) = (F(\pi_1) \circ F(\tilde{\phi}) \circ \eta_{\Z\{S\}})(\beta) = 0  $$

\noindent By step (1), the left outermost square commutes. Hence,
$$(F(\pi_1) \circ F(\tilde{\phi}) \circ \eta_{\Z\{S\}})(\beta) = (F(\pi_1) \circ \eta_{\Z\{R\}} \circ G(\tilde{\phi})(\beta) = 0 $$
By definition $\alpha = G(\tilde{\phi})(\beta)$, hence 
$$F(\pi_1) \circ \eta_{\Z\{R\}}(\beta) = 0 $$
Since $\eta_{\Z\{R\}}(G_I) \subseteq F_I$, the morphism $F(\pi_1) \circ \eta_{\Z\{R\}}$ already factors through $G(\Z\{R\})/G_I$. Thus we have a morphism 
$$C(R) : = G(\Z\{R\})/\langle G_I, \Sigma\rangle \to F(R)$$
Thus, in particular, since $E$ is a pre-Witt functor, there exists a natural transformation $C \to E$. 
\end{proof}

\end{document}
